\documentclass[11pt]{article}
\usepackage[T1]{fontenc}
\usepackage[utf8]{inputenc}
\usepackage{lmodern,microtype,amsmath,amssymb,amsthm,booktabs,tabularx,array,xcolor,fancyhdr}
\usepackage[a4paper,left=24mm,right=24mm,top=22mm,bottom=23mm,headheight=14pt]{geometry}
\usepackage[bookmarksnumbered,colorlinks=true,linkcolor=teal,urlcolor=teal]{hyperref}
\definecolor{ink}{HTML}{142D3F}\definecolor{accent}{HTML}{126C6D}
\usepackage{titlesec}
\titleformat{\section}{\Large\bfseries\color{ink}}{\thesection}{0.6em}{}
\titleformat{\subsection}{\large\bfseries\color{ink}}{\thesubsection}{0.5em}{}
\pagestyle{fancy}\fancyhf{}\fancyhead[L]{\small VMRC \enspace /\enspace Erd\H{o}s problems}\fancyhead[R]{\small Research note}\fancyfoot[C]{\small\thepage}
\setlength{\parindent}{0pt}\setlength{\parskip}{0.38em}
\setcounter{tocdepth}{1}
\hypersetup{pdftitle={Extremal sequences and periodic signs: a VMRC review and proofs on Erdos problems}}
\begin{document}
\begin{center}
{\sffamily\small\color{accent} VERIFIABLE MATHEMATICAL REASONING CORPUS}\par\medskip
{\LARGE\bfseries\color{ink} Extremal sequences and periodic signs}\par
{\Large A corpus review and proofs on Erd\H{o}s problems}\par\medskip
{\small Research manuscript \enspace $\cdot$ \enspace 21 November 2026}
\end{center}
\textbf{Abstract.}
For increasing sequences avoiding sums of consecutive terms, we improve the published Coppersmith--Phillips lower bound and prove that a single infinite sequence attains the finite upper-density envelope. For finite sums of distinct powers we classify all three-base sets at levels one through four and all four-base sets at levels one through three. For the original multiplicative pattern $\{x,2x,3x\}$ we prove uniqueness of the natural-density extremizer up to density zero. We also give an explicit cubic bound for vanishing harmonic sign periods and determine exact extremal profiles for classes of triangle-free Cayley graphs. The first section records the sources and rationale for the corpus ranking. The results below are research arguments with explicit hypotheses and external inputs. Independent publication priority and peer review have not been established.
\tableofcontents
\bigskip
\textbf{Notation.} The prefix VMRC identifies a manuscript in the repository. An Erd\H{o}s problem number identifies a problem in Thomas Bloom's catalogue. Natural density is the limit of $|A\cap[1,N]|/N$ when it exists. Upper and lower natural densities use the corresponding upper and lower limits.
\clearpage
\section{Corpus ranking and transferable methods}
The snapshot contains 225 numbered TeX manuscripts, ending at [225]. The ranking puts the historical strength of the original target, the scope the argument actually covers and the depth of the proof first. Reproducibility helps distinguish closely ranked entries. Closely related manuscripts share one place. These are judgments about the manuscripts. The audit checked the displayed claims and selected proof mechanisms, with a separate exact check of the [055] witness. It was not a line-by-line verification of 225 manuscripts.
\par\smallskip

\noindent\textbf{1.} \href{https://github.com/QEEDDD/Verifiable-Mathematical-Reasoning-Corpus/blob/c2ac34ea4622778187173f41bb1996deee8081b7/\%5B055\%5D\%20ERDOS\%20CONJECTURE.tex}{[055]}, \href{https://github.com/QEEDDD/Verifiable-Mathematical-Reasoning-Corpus/blob/c2ac34ea4622778187173f41bb1996deee8081b7/\%5B056\%5D\%20SQUARE\%20MODULI\%2C\%20SPARSE\%20MASKS\%2C\%20AND\%20SIGN\%20OBSTRUCTIONS\%20FOR\%20VANISHING\%20ERDOS\%20SUMS\%20.tex}{[056]}, \href{https://github.com/QEEDDD/Verifiable-Mathematical-Reasoning-Corpus/blob/c2ac34ea4622778187173f41bb1996deee8081b7/\%5B057\%5D\%20erdos_nonsquare_moduli.tex}{[057]}\quad Erd\H{o}s harmonic nonvanishing. The explicit $q=105^2$ sign table disproves the original universal assertion with $f(0)=0$. The finite character and logarithm conditions can be checked exactly. The non-square branch has also appeared in an earlier external seminar announcement.

\noindent\textbf{2.} \href{https://github.com/QEEDDD/Verifiable-Mathematical-Reasoning-Corpus/blob/c2ac34ea4622778187173f41bb1996deee8081b7/\%5B200\%5D\%20lattice_covering_flows.tex}{[200]}\quad Lattice covering. A bound using discrete surface area gives $\mu(P)\le d/2$ for every non-hollow lattice $d$-polytope and the covering minima of the standard terminal simplex in all dimensions. Compare the original Codenotti--Santos--Schymura paper and arXiv:2601.15173.

\noindent\textbf{3.} \href{https://github.com/QEEDDD/Verifiable-Mathematical-Reasoning-Corpus/blob/c2ac34ea4622778187173f41bb1996deee8081b7/\%5B202\%5D\%20Minc27_R10_Nemeth_Seberry_Shu_Disproof.tex}{[202]}\quad Permanents of prime-order circulants. The manuscript proposes $V(p)\gg p^{4/3}$ against the 1979 $O(p)$ conjecture. Its scope is global. The manuscript labels the disproof a candidate for independent audit.

\noindent\textbf{4.} \href{https://github.com/QEEDDD/Verifiable-Mathematical-Reasoning-Corpus/blob/c2ac34ea4622778187173f41bb1996deee8081b7/\%5B134\%5D\%20Beukers55.tex}{[134]}\quad Beukers--Delaygue supercongruences. Lifted reflection is used to reach $p^{2s}$ for all levels $s$ under the original unit hypothesis. The high-level induction, beyond the initial truncated reflection, is the difficult step.

\noindent\textbf{5.} \href{https://github.com/QEEDDD/Verifiable-Mathematical-Reasoning-Corpus/blob/c2ac34ea4622778187173f41bb1996deee8081b7/\%5B212\%5D\%20randic_extremal_matchings.tex}{[212]}\quad Randi\'c index and matchings. The fractional-matching bound $R(G)\ge\nu_f(G)^{2/3}$ and the sharp limiting coefficient $3\,2^{-1/3}$ address both the exponent and its leading constant.

\noindent\textbf{6.} \href{https://github.com/QEEDDD/Verifiable-Mathematical-Reasoning-Corpus/blob/c2ac34ea4622778187173f41bb1996deee8081b7/\%5B140\%5D\%20ErdosGraham_Density_Parity_Packings.tex}{[140]}, \href{https://github.com/QEEDDD/Verifiable-Mathematical-Reasoning-Corpus/blob/c2ac34ea4622778187173f41bb1996deee8081b7/\%5B185\%5D\%20three_proof_transfers_unified.tex}{[185]}\quad Erd\H{o}s--Graham density. Avoiding $x,2x,3x$ permits maximum natural density $3/4$. The four-point coprime pattern is treated in [185]. The irrationality of the finite limiting constant remains open.

\noindent\textbf{7.} \href{https://github.com/QEEDDD/Verifiable-Mathematical-Reasoning-Corpus/blob/c2ac34ea4622778187173f41bb1996deee8081b7/\%5B101\%5D\%20robin_convex_domain_nonconcavity.tex}{[101]}\quad Robin eigenvalues. Thin convex trapezoids converge to a one-dimensional weighted spectral problem with a strictly positive second variation, furnishing a counterexample to Laugesen\textquotesingle s Conjecture C.

\noindent\textbf{8.} \href{https://github.com/QEEDDD/Verifiable-Mathematical-Reasoning-Corpus/blob/c2ac34ea4622778187173f41bb1996deee8081b7/\%5B196\%5D\%20fractional_green_zeros.tex}{[196]}\quad Fractional Green functions. For $2<\alpha<4$ the manuscript tracks the first loss of positivity and the eventual real-zero count for large parameters. It transfers information between the line and the circle.

\noindent\textbf{9.} \href{https://github.com/QEEDDD/Verifiable-Mathematical-Reasoning-Corpus/blob/c2ac34ea4622778187173f41bb1996deee8081b7/\%5B093\%5D\%20pm1_vanishing_paper.tex}{[093]}, \href{https://github.com/QEEDDD/Verifiable-Mathematical-Reasoning-Corpus/blob/c2ac34ea4622778187173f41bb1996deee8081b7/\%5B107\%5D\%20Dyadic_Saturation_and_Companion_Results_2026_11_02.tex}{[107]}\quad Harmonic sign periods with signs at every residue. Every odd $m\ge3$ occurs as the odd part of a vanishing period. The set of all admissible periods is still not classified.

\noindent\textbf{10.} \href{https://github.com/QEEDDD/Verifiable-Mathematical-Reasoning-Corpus/blob/c2ac34ea4622778187173f41bb1996deee8081b7/\%5B218\%5D\%20De_Bruijn_Discrepancy_Complete.tex}{[218]}\quad De Bruijn discrepancy. The Prefer-opposite discrepancy is $n^2/2+O(n)$, resolving Gabric--Sawada 3.2. The Prefer-same clause of 3.1 is covered. The Lex-comp clause remains.

\subsection*{Methods that transfer beyond the ranked statements}
\begin{tabularx}{\linewidth}{@{}p{.23\linewidth} X X@{}}
\toprule
Corpus source & Transferable operation & Test on the next problem\\\midrule
{[055], [093], [107]} & Replace a global harmonic cancellation by a finite identity on valuation layers. & Supply exact weights and control amplitude. A balanced formal identity alone need not define signs.\\
{[140], [185], [193]} & Charge each forbidden multiplicative pattern to a missing integer. Follow the cases of zero slack. & In a stability question, count missing and extra integers in the same initial segment and track dilation.\\
{[203]} & Pass to a quotient on which the cost of a cut can be realized, and lift the partition. & Count odd cycles to obtain the reverse inequality, or exhibit a sign-valued lowest eigenvector.\\
{[200]} & Average values at jumps of a periodic function, preserving the linear drift between jumps. & Needs an exact period and control at every jump. It is suited to covering inequalities.\\
{[134]} & Reflect paired residue classes while retaining the term that is lost in the first reduction. & For a higher-level congruence, prove the correction gains the required power of $p$ on iteration.\\
{[108], [193]} & Use a finite certificate only after proving that it propagates. & For sums of powers the certificate is a whole interval. In the consecutive-sum problem it is a finite residue table with a bound on later deletions.\\
\bottomrule
\end{tabularx}
\par
In particular the prime-lift existence question in [093] was answered by [107]. The $(2,3,5)$ multiplicative density $61/72$ was already in [185], and the diagonal weighted corner problem from [140] was already treated in [195]. The theorems below use the corresponding later statements as baselines.

\clearpage

\section{Consecutive sums and the finite-to-infinite extremal bound}

Call an increasing sequence of positive integers \emph{admissible} if the sum of at least two consecutive terms is never a term. Let $f(N)$ be the largest size of an admissible subset of $\{1,\dots,N\}$. Coppersmith and Phillips, \emph{SIAM J. Discrete Math.} 9 (1996), 173--177, Theorem 2.1, proved $f(N)\ge (13/24)N-O(1)$. Their Open Question 1 asks for $\liminf f(N)/N$. Freud had already answered the original Erd\H{o}s problem 867 negatively in 1993.

\paragraph{Theorem.} For every $N$,
\[
f(N)\ge \frac{703}{1296}N-O(1).
\]
A single infinite admissible set attains upper asymptotic density $\limsup_N f(N)/N$, which is therefore at least $703/1296$. The exact extremal coefficient and the questions about lower and logarithmic density in Erd\H{o}s problem 839 remain undetermined.

\subsection{The construction}
For a normalized open interval $J$ and modulus $q$, choose the integers $t$ such that $t/N\in J$ and $t\bmod q$ lies in the given set. Omit interval endpoints. A bar denotes complement modulo the modulus of that row. Use the following baseline from Coppersmith and Phillips, Table 1, with no entries below $1/5$. Set
\[
\begin{aligned}
C&=\{0,2,6,8,9,12,13,16,19,20,23,24,26,30\},\\
D&=\{2,8,14,17,21,25,29,35,39,43,47,50,56,62\}.
\end{aligned}
\]
\begin{center}\small
\begin{tabular}{lll}
Normalized interval & Modulus & Residues kept \\ \hline
$(1/5,2/9)$&16&$0,1,2,5,6,8,10,11,14,15$\\
$(2/9,1/4)$&1&all\\
$(1/4,8/27)$&1&none\\
$(8/27,1/3)$&4&$1,2,3$\\
$(1/3,3/8)$&3&$1,2$\\
$(3/8,4/9)$&32&$C$\\
$(4/9,1/2)$&2&$0$\\
$(1/2,16/27)$&1&all\\
$(16/27,2/3)$&8&$1,2,4,6,7$\\
$(2/3,3/4)$&3&$1,2$\\
$(3/4,8/9)$&64&$\overline D$\\
$(8/9,1)$&4&$0,1,3$\\
\end{tabular}\end{center}
Now overwrite three bands. Put $I=(16/81,1/5)$ and $kI=\{kx:x\in I\}$. Define
\[
\begin{aligned}
X&=\{1,4,6,9,14\},\\
U&=\{0,3,6,9,11,16,17,22,24,28\},\\
V&=\{3,9,15,20,27,33,39,46,52,60\}.
\end{aligned}
\]
Use $X\bmod16$ on $I$, $U\bmod32$ on $2I$, and $\overline V\bmod64$ on $4I$. Omit all patch endpoints. The original density contribution of these bands was $|I|(2\cdot14/32+4\cdot50/64)=4|I|$. The last two modified bands still contribute $|I|(2\cdot10/32+4\cdot54/64)=4|I|$. Thus the first band gives the full gain
\[
\frac{5}{16}|I|=\frac{5}{16}\left(\frac15-\frac{16}{81}\right)=\frac1{1296}.
\]
Each residue count differs from its interval-length estimate by a bounded quantity. Summing the twelve original length-density products gives $13/24$, so the raw construction has $(703/1296)N+O(1)$ terms.

\subsection{Verification}
The forbidden four-term sums from $I$ already occur among the forbidden two-term sums from $2I$. The following table verifies this and all remaining periodic cases. Use the residue sets
\[E=\{1,3,7,11,14,18,21,25,29,31\},\qquad F=\{2,8,14,21,29,35,43,50,56,62\}.\]
For each $k\ge2$, list intervals on which the source pattern and the pattern at $kx$ are both nonempty with $kx<1$. The right column gives the residues of sums of $k$ consecutive entries in the periodic source pattern, modulo the target modulus. The target pattern excludes every listed residue. Each row is checked by listing the accepted source residues over one common period of the source and target moduli, appending $k-1$ successors, and summing each length-$k$ window.
\begin{center}\small
\begin{tabular}{rlll}
$k$ & Source interval & Target modulus & Sum residues \\ \hline
2 & $I$ &32&$5,10,15,23,31$\\
2 & $(1/5,2/9)$ &32&$E$\\
2 & $(2/9,1/4)$ &2&$1$\\
2 & $(8/27,1/3)$ &8&$0,3,5$\\
2 & $(1/3,3/8)$ &3&$0$\\
2 & $(3/8,32/81)$ and $(2/5,4/9)$ &64&$D$\\
2 & $2I$ &64&$V$\\
2 & $(4/9,1/2)$ &4&$2$\\
3 & $I$ and $(1/5,2/9)$ &8&$0,3,5$\\
3 & $(2/9,1/4)$ &3&$0$\\
3 & $(8/27,1/3)$ &4&$2$\\
4 & $I$ &64&$9,20,33,46,60$\\
4 & $(1/5,2/9)$ &64&$F$\\
4 & $(2/9,1/4)$ &4&$2$\\
5 & $I$ &4&$2$\\
\end{tabular}\end{center}

There are no relevant blocks of length at least six, since $6(16/81)>1$. For each $2\le k\le5$, partition $[0,1]$ at every interval endpoint $e$ of the modified construction and at every $e/k$. If the terms of a raw block lie in one open cell, their normalized mean $x$ lies in that cell. The target $kx$ lies in a single target band. The raw summands occur consecutively in the periodic source pattern. The table therefore excludes their sum, or the target is empty or outside the range. At a cut, only $O(1)$ blocks of length at most five can straddle it or contain an entry exactly at the cut. The number of cuts is fixed. Thus only $O(1)$ raw selected targets can be represented by raw consecutive sums.

\paragraph{Deletion lemma.}
Suppose $T\subseteq[\delta N,N]$ consists of distinct integers with $0<\delta<1$, and at most $h$ of its terms are sums of at least two consecutive terms of $T$. Scan $T$ in increasing order, rejecting a term if it is a sum of at least two consecutive previously accepted terms. The resulting sequence is admissible and the number of rejections is at most
\[
 h\sum_{j=0}^{L}(L^2)^j,\qquad L=\lceil1/\delta\rceil.
\]
Indeed, choose a witness block for each rejected $t$. If its summands were consecutive in $T$, declare $t$ a root. Otherwise a rejected $z$ lies strictly between the first and last summands. Choose one such $z$ as parent. Since the largest summand exceeds $z$ and another is at least $\delta N$, one has $t>z+\delta N$. Every witness remains consecutive in the final accepted list. A given parent lies in one gap of this list. At most $L^2$ blocks of length at most $L$ span the gap, and each block determines its sum uniquely. Each parent therefore has at most $L^2$ children. The value increase bounds the depth by $L$. Taking $\delta=16/81$ and $h=O(1)$ proves the finite bound, including all deletions.

\subsection{One infinite sequence attaining the finite upper-density envelope}
Choose a finite admissible prefix $P$ with $s$ terms and fix $0<\delta<1$. Take $N$ large enough that $\delta N>\sum_{p\in P}p$. Discard entries below $\delta N$ from an extremal admissible set in $[1,N]$, leaving an admissible tail $T$ of size at least $f(N)-\delta N-1$. Apply the same scan with $P$ already accepted. A sum made entirely from $P$ cannot hit $T$. A witness crossing into the new accepted list uses one of $s$ old suffixes and one of at most $L$ nonempty new prefixes, where $L=\lceil1/\delta\rceil$. At most $sL$ distinct rejections can have this kind of witness. Every remaining rejection spans an earlier deletion in $T$, because $T$ was admissible. The gap argument gives at most $sL\sum_{j=0}^{L}(L^2)^j$ rejections altogether, independent of $N$.

Set $d=\limsup_N f(N)/N$ and choose $\delta_j\downarrow0$. At stage $j$, keep the previous prefix. Choose $N_j$ arbitrarily large with $f(N_j)/N_j>d-\delta_j$, with $\delta_jN_j$ greater than both the old prefix sum and $N_{j-1}$, and with the bounded rejection count plus one less than $\delta_jN_j$. Append the accepted tail. The infinite union is admissible since every possible violation is finite. Its density along the $N_j$ is at least $d-3\delta_j$. Conversely every prefix of any infinite admissible sequence has size at most $f(N)$. This proves attainment.

\paragraph{Sources.}
D. Coppersmith and S. Phillips, \emph{On a Question of Erd\H{o}s on Subsequence Sums}, SIAM Journal on Discrete Mathematics 9 (1996), 173--177, DOI 10.1137/S0895480193244139.
R. Freud, \emph{Adding numbers}, James Cook Mathematical Notes 6 (1993), 6199--6202. See also the public problem records for Erd\H{o}s numbers 839 and 867. The comparison establishes a strict improvement over the published 1996 coefficient. Publication priority for the modified construction has not been established.



\subsection{Three further modifications and a coefficient above \(0.543\)}
The same construction admits the stronger bound
\[
f(N)\ge \frac{3040369}{5598720}N-O(1).
\]
The infinite transfer gives the same lower bound for attainable upper density. For a periodic set $A$ on the integers, let $S_k(A)$ be the periodic set of sums of $k$ successive terms. If $A$ has period $q$ and $r$ entries in one period, $S_k(A)$ has period $kq$ with $r$ entries per such period. This follows from strict increase of the sliding sums and translation by $q$. We describe additional replacements on dilates of three intervals.
\[
J=(3/16,16/81),\qquad H=(8/45,3/16),\qquad K=(1/6,128/729).
\]
Write $[R]_q=\{n:n\bmod q\in R\}$ and $Y=[\{29,33,37,41\}]_{64}$. In the following table a blank means to retain the preceding construction. A superscript $c$ means complement in the integers. Set $A=[\{2\}]_4$, $B=[\{2,10\}]_{12}$, $C'=[\{88,152\}]_{192}$, $U'=[\{3,5\}]_8$, and $W=[\{0\}]_8$.
\begin{center}\small
\begin{tabular}{c|ccccc}
$L$ & $L$ & $3L/2$ & $2L$ & $3L$ & $4L$\\ \hline
$J$&$A$&$Y$&$U'$&$(S_3(A)\cup S_2(Y))^c$&$S_2(U')^c$\\
$H$&$B$&$Y$&&$(S_3(B)\cup S_2(Y))^c$&\\
$K$&$C'$&$W$&&$S_2(W)^c$&
\end{tabular}\end{center}
All overwritten bands in this table and in the original three replacements are disjoint, apart from endpoints, which are omitted. The least selected entry is now greater than $N/6$, so lengths two through five still suffice.

For $J$, the pair sums from $A$ are $0\bmod8$, which $U'$ excludes. The four-term sums from $A$ are $0\bmod16$ and $S_2(U')$ consists of $0,8\bmod16$. The five-term sums are $2\bmod4$, excluded on $5J$. For $H$ and $K$, each successive pair of the source terms sums to $0\bmod3$ (on both periodic patterns the nonzero residues alternate). Hence the pair and four-term sums miss the unchanged bands $2H,4H,2K,4K$, which retain $1,2\bmod3$. Five-term sums from $B$ are $2\bmod4$, excluded on $5H$. The terms of $C'$ are all $24\bmod64$, so five-term sums from $C'$ are $56\bmod64$, excluded on $5K$ by $D$. Its triple sums are $8\bmod16$, exactly the residue class $S_2(W)$.

The triple sums from $Y$ have residues $35,39,43,47\bmod64$, excluded on $9J/2$ and $9H/2$ by $D$. The triples from $W$ are $0\bmod8$, excluded on $9K/2\subseteq4J$ by $S_2(U')$. On $3J,3H,3K$, the only nonempty sources with block length at least two are $L$ (triples) and $3L/2$ (pairs). On $2J$ the only such source is $J$ (pairs). On $4J$ the new terms from $J$ and $2J$ yield the forbidden residues just verified. Triples landing in $4J$ have normalized source in $(1/4,64/243)$, where a new source from $3K/2$ may occur. This is $W$, treated above. For lengths at least five the source for $4J$ is below $1/6$. The new bands $L$ and $3L/2$ themselves have no incoming sums, as their upper endpoints are less than $1/3$. All sums of length at least four from $3L/2$ already exceed $N$, as do all sums of length at least two from $3L$ or $4J$. These checks and the unchanged cases verify all homogeneous blocks. The boundary and deletion arguments above give a loss of $O(1)$.

For the count, the baseline contribution of the $2J,4J$ bands remains $4|J|$, since selecting $U'$ on $2J$ costs exactly the same count in its two-term sums on $4J$. Each $3L$ band was full, and each $L,3L/2$ band was empty. The densities in the new $L$ and $3L/2$ bands exactly pay for their forbidden triple and pair sums on $3L$, with a gain equal to the overlap. On $J$, $S_3(A)$ is $6\bmod12$. On $H$, $S_3(B)$ is $10,26\bmod36$ (both $2\bmod4$). The four pair sums of $Y$ are $6,62,70,78\bmod128$, all $2\bmod4$. By the Chinese remainder theorem, the overlap densities on $3J$ and $3H$ are respectively $4/384$ and $8/1152$. On $3K$ every triple sum of $C'$ is already excluded by $S_2(W)$, so the whole density contribution from $C'$ is gained. The three gains are
\[
3|J|\frac{4}{384}=\frac{13}{41472},\quad
3|H|\frac{8}{1152}=\frac7{34560},\quad
|K|\frac2{192}=\frac{13}{139968}.
\]
Adding these to $703/1296$ gives $3040369/5598720$. The appendix checks all 35 active interval and block-length combinations. At $N=10^6$ the endpoint convention above gives 543,044 raw entries and 543,010 after greedy deletion. A separate check of every remaining sliding sum of lengths two through five finds no violation.

\clearpage

\section{A sharp threshold and stability for multiplicative forbidden patterns}
Let $s\ge2$, and let $2\le a_1<\cdots<a_s<b$ be pairwise coprime integers (including $b$). Consider subsets $A$ of the positive integers avoiding every configuration
\[
 \{x,a_1x,\ldots,a_sx,bx\}\qquad(x\ge1).
\]
Write $\nu_c(n)=\max\{j:c^j\mid n\}$ for $c\ge2$, and put
\[
 B=\{n:\nu_b(n)\text{ is even}\},\quad D_0=\frac{b}{b+1},\quad
 Q=1+\sum_{i=1}^s a_i,\quad q=Q/b.
\]
The symbol $d$ denotes natural density when it exists, and $\overline d$ denotes upper natural density. The following strengthens the sufficient condition in corpus [140]. Corpus [185] already determines the optimum for $s=2$ at every parameter value. The threshold test below extends to all $s$, and the quantitative stability statement concerns all competing sets.
\par\medskip\noindent\textbf{Theorem.}
The set $B$ has maximum natural density among the admissible sets possessing natural density if and only if $b\ge Q$. When $b>Q$, every such $A$ satisfies the sharp estimate
\[
 \overline d(A\mathbin{\triangle}B)\le
 \frac{b+Q}{b-Q}\bigl(D_0-d(A)\bigr).
\]
In particular, every extremizer agrees with $B$ outside a set of natural density zero. If $b=Q$, two extremizers differ on a set of positive natural density. If $b<Q$, an explicit admissible set has density
\[
 D_0+\frac{Q-b}{(b+1)\prod_{i=1}^s(a_i+1)}.
\]
This last expression is a lower bound for the optimum in the region $b<Q$.
\par\medskip\noindent\emph{Proof.}
Multiplication by $b$ changes $\nu_b$ by one, so $B$ is admissible.
Let $U=B\setminus A$ and $V=A\setminus B$, and use $S(X)=|S\cap[1,X]|$.
For every $y\in V$, the value $x=y/b$ is integral with even $b$-valuation. Each of $x,a_1x,\ldots,a_sx$ also has even $b$-valuation. At least one is in $U$, because $bx=y\in A$ and $A$ avoids the full configuration. All the maps $y\mapsto y/b$ and $y\mapsto a_iy/b$ are injective. Hence
\[
 V(X)\le U(X/b)+\sum_{i=1}^s U(a_iX/b).
\]
Suppose $d(A)$ exists and write $\epsilon=D_0-d(A)$. The identity
$U(X)-V(X)=B(X)-A(X)=\epsilon X+o(X)$ follows from $d(B)=D_0$.
If $u=\limsup U(X)/X$, taking upper limits in the counting inequality gives
$(1-q)u\le\epsilon$. In particular $\epsilon\ge0$ whenever $q\le1$.
If $q<1$, then $\overline d(A\triangle B)=2u-\epsilon\le(1+q)\epsilon/(1-q)$.

For sharpness, define $H(n)=|\{i:\nu_{a_i}(n)\text{ is even}\}|$ and
\[
 T=\{n:\nu_b(n)\text{ even},\ H(n)=0\}
 \ \cup\ \{n:\nu_b(n)\text{ odd},\ H(n)\ge2\},
 \qquad A_* = \mathbb N\setminus T.
\]
We check that $T$ meets each configuration. When $H(x)=0$ or $H(x)\ge2$,
one of $x,bx$ lies in $T$. For $H(x)=1$ and $\nu_b(x)$ even, multiply by the unique $a_i$ with even valuation, obtaining $H=0$. For $H(x)=1$ and $\nu_b(x)$ odd, multiply by an $a_i$ with odd valuation, obtaining $H=2$; such an index exists since $s\ge2$. Pairwise coprimality ensures that each multiplication toggles exactly its own parity.

For any individual base $c$, the sets where $\nu_c$ is odd and even have densities $1/(c+1)$ and $c/(c+1)$, respectively. The parities of our pairwise coprime bases are independent in natural density. Indeed, exact truncated valuation classes factor by the Chinese remainder theorem, and the integers divisible by a $K$th power of one of the bases have upper density at most $b^{-K}+\sum_i a_i^{-K}$. Let $P=\prod_i(a_i+1)^{-1}$. The densities of $H=0$ and $H=1$ are $P$ and $(Q-1)P$, respectively. For $A_*$ we therefore have
\[
 d(B\setminus A_*)=\frac{bP}{b+1},\qquad
 d(A_*\setminus B)=\frac{QP}{b+1}.
\]
These equalities give the advertised density, prove sharpness of the constant in the strict region, and exhibit both the positive-density ties at equality and the failure of optimality below the threshold.\hfill$\square$

\subsection*{Rigidity for the original three-point configuration}
With only one multiplier besides $b$, the critical case has a different outcome. Corpus [193] already proves uniqueness and sharp H\"older stability on a single weighted exponent lattice. The following counts integers in ordinary initial intervals. Their behaviour may differ across the infinitely many multiplicative fibres, so a result on one fibre alone does not establish this formulation.
\par\medskip\noindent\textbf{Theorem.}
Let $a\ge2$ and $b=a+1$. If $A\subseteq\mathbb N$ avoids $\{x,ax,bx\}$ and $d(A)=b/(b+1)$, then $d(A\triangle B)=0$, where $B=\{n:\nu_b(n)\text{ is even}\}$. In particular every natural-density maximizer for the original configuration $\{x,2x,3x\}$ agrees, outside a set of natural density zero, with the integers of even $3$-adic valuation.
\par\medskip\noindent\emph{Proof.}
Use $U,V$ as above. Since $d(A)=d(B)$, $U(X)-V(X)=o(X)$. Write $u=\overline d(U)$ and choose integers $N_j\to\infty$ with $U(N_j)/N_j\to u$. The counting inequality is now
\[
 V(X)\le U(X/b)+U(aX/b).
\]
After division by $X$, the two coefficients of the rescaled densities on the right are $1/b$ and $a/b$, with sum one. Each rescaled density has upper limit at most $u$. Thus when $X=N_j$, both tend to $u$, and the slack
\[
 L(N_j)=U(N_j/b)+U(aN_j/b)-V(N_j)
\]
is $o(N_j)$. Applying the same argument at $N_j/b$, and then iterating, shows $U(N_j/b^M)/(N_j/b^M)\to u$ for every fixed $M\ge1$.

Call $n\in U$ good if $a\mid n$ and both $bn$ and $bn/a$ belong to $V$. There are at most $2L(N_j)$ non-good elements of $U$ up to $N_j/b$. To see this, count two lists of potential images: each $n\in U$ with $n\le N_j/b$ supplies a place for $bn$, and each $n\in U$ with $n\le aN_j/b$ supplies a place for $bn/a$. A place is unused if its image fails to be an integer of $V$. Every member of $V$ up to $N_j$ has at least one place, and the number of unused places is at most $L(N_j)$. Non-good elements in the smaller range give unused places, including when $a\nmid n$, since $(a,b)=1$.

For a good $n$ the integers $bn/a$ and $bn$ lie in $A$. Avoidance forces $b^2n/a\notin A$. Its $b$-valuation is even, so $b^2n/a\in U$. Set $F(n)=b^2n/a$ on good $n$. Fix $h\ge1$ and choose $M$ with $(b^2/a)^h b^{-M}\le1/b$. Starting below $N_j/b^M$, the first $h$ steps stay below $N_j/b$ as long as the iterates are good. For each fixed step, the map into that step is injective. Hence at most $2hL(N_j)=o(N_j)$ initial elements of $U$ encounter a non-good element in these steps. Every other initial element is divisible by $a^h$, because $F^h(n)=b^{2h}n/a^h$ is integral and $(a,b)=1$. Therefore
\[
 U(N_j/b^M)\le \frac{N_j}{b^M a^h}+o(N_j).
\]
Letting $j\to\infty$ gives $u\le a^{-h}$. Letting $h\to\infty$ gives $u=0$, and $V(X)=U(X)+o(X)$ finishes the proof.\hfill$\square$

\clearpage

\section{A full three-base classification at the first four positive levels}

For a finite set $D$ of distinct integers at least three, let $P_k(D)$
be the set of finite sums of the terms $d^j$, $d\in D$, $j\geq k$,
using each indexed term at most once. Write
$\sigma(D)=\sum_{d\in D}1/(d-1)$. Burr, Erd\H{o}s, Graham and Li
conjectured that $P_k(D)$ contains all sufficiently large integers
precisely when $\sigma(D)\geq1$ and $\gcd(D)=1$
(\emph{Acta Arith.} 77 (1996), 133--138).
Their concluding remarks single out $D=\{3,4,7\}$ and report the
largest missing integer $581$ at $k=1$. Their statement there refers
to positive exponents starting at one.

\paragraph{Theorem.}
For each $k\in\{1,2,3,4\}$ and every three-element set $D$ of
integers at least three, $P_k(D)$ is cofinite if and only if
$D=\{3,4,c\}$ for $c\in\{5,6,7\}$. The exact largest missing
integers are as follows.
\[
\begin{array}{c|r|r|r}
 k&c=5&c=6&c=7\\ \hline
 1&79&986&581\\
 2&77\,613&242\,113&3\,982\,888\\
 3&4\,330\,731&58\,941\,162&166\,025\,260\\
 4&1\,075\,364\,603&14\,059\,625\,243&14\,961\,204\,826
\end{array}
\]
The two strict-inequality columns were previously reported in the
2026 computational notes at \texttt{mrly.net/research/claims/erdos/}.
The verification below also computes those columns. The three
higher levels of the critical column are not listed in the sources
examined. Historical priority is unverified. The theorem has no
restriction on the sizes of the three bases. The case of all positive
levels, even for the triple $\{3,4,7\}$, remains outside its scope.

\paragraph{Necessity.}
If $\sigma(D)<1$, there are arbitrarily large real $x$ for which, for
each $d$, the least power $u_d=d^j\geq x$ has $u_d/x\to1$ along these
scales. Indeed, simultaneous Dirichlet approximation gives unbounded
integers $Q$ for which the distances of $Q/\log d$ to the nearest
integers all tend to zero. Take $x$ to be the smallest of the
corresponding nearest powers to $e^Q$ on the logarithmic scale.
Their ratios to $x$ tend to one, so eventually each preceding power
lies below $x$. The total $S$ of all terms below $x$ satisfies
$S/x\to\sigma(D)$ by the geometric series formula. Every integer
strictly between $S$ and $x$ is missing. This gives the necessity
credited by the original paper to Pomerance. If $3\le a<b<c$
and $a\ge4$, then $\sigma\le1/3+1/4+1/5<1$. If $a=3$ and $b\ge5$,
then $\sigma\le1/2+1/4+1/5<1$. For $(a,b)=(3,4)$ the requirement
is $c\le7$. This leaves the three displayed triples.

\paragraph{Sufficiency and sharpness.}
Order the terms in increasing order. If a prefix of sum $S$ realizes
every integer in $[h,S-h]$, adjoining $t$ extends that interval to
$[h,S+t-h]$ whenever $S-2h+1\ge t$. This follows by merging it
with its translate by $t$. If also $h-1$ is missing from this prefix
and $h\le t$, then $h-1$ remains missing.

Let $u_d$ be the least unprocessed power for base $d$ and put
$t=\min_d u_d$ and $C=\sum_d d^k/(d-1)$. At every step,
\[
 S-t=(\sigma(D)-1)t+\sum_{d\in D}\frac{u_d-t}{d-1}-C .
\]
For all three triples $S-t\ge |u_3-u_4|/3-C$.
Write $u_4=2^a$ and $u_3=3^b$. Ellison's explicit estimate gives
$|2^a-3^b|>2^a e^{-a/10}>2^{a/2}$ for all $a>27$.
See W.~J.~Ellison, \emph{On a theorem of S.~Sivasankaranarayana Pillai}
(1971), and M.~Waldschmidt, \emph{Perfect powers: Pillai's works and
their developments}, \S2.2.2. For all twelve cases $C\le526$
and $h\le14\,961\,204\,827$. Consequently when $t\ge2^{100}$,
we have $S-t>2^{50}/3-526>2h-1$. All later merges are guaranteed.

The following integer certificate records the prefix length $n$,
prefix sum $S$, next term $t$, number $J$ of further terms below
$2^{100}$, and the minimum $M$ of $S-2h+1-t$ over those $J$ steps.
Here $h$ is one plus the entry of the first table.
\[
\begin{array}{rr|r|r|r|r|r}
 c&k&n&S&t&J&M\\ \hline
5&1&11&602&256&144&74\\
5&2&23&450742&262144&129&33371\\
5&3&29&24973824&14348907&120&1963454\\
5&4&39&8187710991&4294967296&107&1742014488\\
6&1&16&6197&4096&134&128\\
6&2&26&2531169&1594323&121&452619\\
6&3&35&624183519&362797056&109&143504138\\
6&4&44&154370390829&78364164096&97&23880576321\\
7&1&15&5042&2401&132&553\\
7&2&31&33841349&16777216&113&2299181\\
7&3&37&2342340864&1162261467&104&848028876\\
7&4&43&154845726270&94143178827&95&30780137790
\end{array}
\]
Here is self-contained Python code for the finite proof. It keeps
only the sorted positive holes in the lower half of each prefix.
For the old total $S$, a point $x\le S$ in the new lower
half can be missing only if $x$ was an old lower hole or $S-x$ was
one. Among these candidates $x$ remains missing exactly when $x-t$
is negative or an old hole, since $x-t\le S/2$. When $x>S$ in
the new lower half, $x<t$ and it is automatically missing. These
observations prove the update used in the code. Reflection in $S/2$
supplies the upper half of the certified central interval.
{\small\begin{verbatim}
from array import array

def holes_after(holes, s, t):
    hi = (s + t) // 2
    out = array("Q")
    def retain(candidates):
        j = 0
        for x in candidates:
            v = x - t
            while j < len(holes) and holes[j] < v:
                j += 1
            if v < 0 or (j < len(holes) and holes[j] == v):
                out.append(x)
    retain(holes)
    def reflected():
        for z in reversed(holes):
            x = s - z
            if x > hi:
                break
            if x > s // 2:
                yield x
    retain(reflected())
    out.extend(range(s + 1, hi + 1))
    return out

def certificate(c, k, n):
    bases = (3, 4, c)
    u = [d**k for d in bases]
    holes, s = array("Q"), 0
    for _ in range(n):
        t = min(u)
        i = u.index(t)
        holes = holes_after(holes, s, t)
        s += t
        u[i] *= bases[i]
    h = holes[-1] + 1 if holes else 1
    assert h <= min(u) and 2*h <= s
    seed = (s, min(u))
    margins = []
    while min(u) < 2**100:
        t = min(u)
        i = u.index(t)
        margins.append(s - 2*h + 1 - t)
        s += t
        u[i] *= bases[i]
    assert min(margins) >= 0
    return h-1, seed, len(margins), min(margins)

for c, ns in [(5,(11,23,29,39)), (6,(16,26,35,44)),
              (7,(15,31,37,43))]:
    for k,n in enumerate(ns,1):
        print(c, k, certificate(c,k,n))
\end{verbatim}}

Thus the finite checks certify an interval and its largest lower hole,
and the analytic estimate makes the asserted coverage infinite.



\section{The four-base classification through level three}

\paragraph{Theorem.}
Let $D$ be any four distinct integers at least three. For each
$k=1,2,3$, $P_k(D)$ is cofinite if and only if $\sigma(D)\ge1$.
Every qualifying four-base set automatically has gcd one.

If $D$ contains one of the three successful triples from the preceding
section, its conclusion persists upon adding the fourth base. For the
other qualifying quadruples, write $D=\{a<b<c<d\}$. We must have
$a=3$ since the weights of $4,5,6,7$ sum to less than one. We must have
$b\le6$ since the weights of $3,7,8,9$ also sum to less than one.
For $b=4$ the absence of a successful triple forces $c\ge8$.
The deficit of the first three weights is then at least $1/42$, so
$d\le43$. For $b=5$, $c\ge6$ gives a deficit of at least $1/20$
and $d\le21$. For $b=6$, $c\ge7$ gives a deficit of at least
$2/15$, hence $d\le8$, forcing $(c,d)=(7,8)$. This also proves
the gcd assertion. Enumeration now gives precisely $91$ quadruples,
distributed as $67,23,1$ over $b=4,5,6$. Our enumeration attaches
distinct indices to equal powers of distinct bases, as required by
the sum of the four digit-restricted sets.

\paragraph{Uniform tail bound.}
For all the quadruples, put $u_3=3^i$ and $u_b=b^j$ for the first
unprocessed powers, and $t=\min_{e\in D}u_e$. The previous identity
gives $S-t\ge |3^i-b^j|/5-C$. We show
$|3^i-b^j|>10^{30}$ when $t\ge2^{100}$.
First sweep the two increasing power streams of bases 3 and $b$
from their first terms at least $2^{100}$. At each step compare the
two current values, then advance the smaller. The successive states
include every pair of least powers above a common threshold, as
occurs at our merges. Up to the first state at which the smaller
term is at least $2^{50000}$, the minima of the absolute differences
are, in order for $b=4,5,6$,
\[
\begin{array}{r|r|r}
 b &\text{comparisons}&\text{minimum difference}\\ \hline
 4&56433&1636918580620405121328963732223\\
 5&52973&2250658065788289002311145051344\\
 6&50787&1205768063818556806684708093185
\end{array}
\]
Even if a power matching another indexed base has just been processed,
the two selected streams are at one of the same successive states.
For a state beyond this sweep let $x=\min(3^i,b^j)\ge2^{50000}$
and $L=\log x$. Since the two streams are advanced in order,
$3^i\le3x$ and $b^j\le bx$. Use Laurent's two-logarithm estimate
in the positive rational form of Theorem~4 of Bennett, Bugeaud and
Mignotte, \emph{Perfect powers with few binary digits and related
Diophantine problems, II}, Math. Proc. Cambridge Philos. Soc. (2012),
or directly Corollary~2 of M.~Laurent, \emph{Linear forms in two
logarithms and interpolation determinants II}, Acta Arith. 133 (2008).
Specifically, with $B'=i/\log b+j/\log3$,
\[
 \log |i\log3-j\log b|
 \ge -25.2\log3\log b\,
       \max\{\log B'+0.38,10\}^2 .
\]
Here $B'\le4L$, $25.2\log3\log b<50$, and $L>34000$.
At $L=34000$, $50(\log(4L)+0.38)^2<L/2$, and the ratio of
the left side to $L$ decreases beyond that point. Thus
$|3^i-b^j|\ge x|i\log3-j\log b|>\exp(L/2)>10^{30}$.
For these quadruples and levels $C\le4\cdot43^3$ and the largest
certified $h$ is $1\,156\,630\,106$. Hence every merge after the
$2^{100}$ threshold is valid.

\paragraph{Finite computation.}
The columns below give the largest prefix length used, the sum and
maximum of the 91 exact largest exceptions $F=h-1$, and the minimum
finite merge margin among the selected certificates.
\[
\begin{array}{r|r|r|r|r|r}
 k &\text{count}&\max n &\sum F&\max F&\min M\\ \hline
 1&91&18&12466&879&0\\
 2&91&34&14346879&3948300&117\\
 3&91&43&4960642426&1156630105&7807
\end{array}
\]
The following code uses the function \texttt{holes\_after} of the
preceding section. It regenerates each certificate, including zeros
of the merge margin, where the two integer intervals are adjacent.
{\small\begin{verbatim}
from fractions import Fraction
from itertools import combinations

def bound_check(b):
    x, y = 3, b
    while x < 2**100:
        x *= 3
    while y < 2**100:
        y *= b
    m, steps = None, 0
    while min(x,y) < 2**50000:
        v = abs(x-y)
        m = v if m is None else min(m,v)
        steps += 1
        if x < y:
            x *= 3
        else:
            y *= b
    assert m > 10**30
    return steps, m

def find_seed(D, k):
    u = [d**k for d in D]
    holes, s = array("Q"), 0
    for n in range(44):
        t = min(u)
        h = holes[-1]+1 if holes else 1
        if n >= 6 and h < t and s-2*h+1 >= t:
            z, total, margin = u.copy(), s, []
            while min(z) < 2**100:
                y = min(z)
                j = z.index(y)
                margin.append(total - 2*h + 1 - y)
                if margin[-1] < 0:
                    break
                total += y
                z[j] *= D[j]
            if min(margin) >= 0:
                return n, h-1, min(margin)
        i = u.index(t)
        holes = holes_after(holes,s,t)
        s += t
        u[i] *= D[i]
    raise ArithmeticError("no certificate at this search depth")

def weight(D):
    return sum(Fraction(1,d-1) for d in D)

cases = [D for D in combinations(range(3,44),4)
         if weight(D[:3]) < 1 <= weight(D)]
assert len(cases) == 91
for b in (4,5,6):
    print(b, bound_check(b))
for k in (1,2,3):
    rows = [find_seed(D,k) for D in cases]
    indices, largest, margins = zip(*rows)
    print(k,len(rows),max(indices),sum(largest),
          max(largest),min(margins))
\end{verbatim}}
The negative direction follows from the same missing-interval argument
as for triples. The classification includes four-base sets of
arbitrarily large maximum base via the containment step.



\section{A cubic bound for vanishing harmonic sign periods}

Let $\mathcal D$ be the set of integers $N$ for which an $N$-periodic
$f:\mathbb Z\to\{-1,1\}$ satisfies $\sum_{n=1}^N f(n)=0$ and
$\sum_{n\ge1}f(n)/n=0$. The signs include the residue $0$.
Murty and Pathak discuss this variant of the Erd\H{o}s problem in
Section~4.4 of \emph{The Okada space and vanishing of $L(1,f)$}.
Corpus item [107] showed that for every odd $m\ge3$ there exists
$a\le5\lceil\log_2 m\rceil+10$ with $2^a m\in\mathcal D$. Its
construction at $m=127$ gave $a=41$.

\paragraph{Signed-digit transfer [107].}
Write $U_a=(\mathbb Z/2^a\mathbb Z)^\times$ and
$w(x)=2^{v_2(x)}$ for $x\ne0$.
If $H:U_a\to\mathbb Z$ is antisymmetric and
$H(u)\equiv H(v)\pmod{2^j}$ whenever $u\equiv v\pmod{2^j}$ for
$1\le j\le a$, put $G(u)=H(mu)-mH(u)$. The proposition
\emph{Transfer to an arbitrary odd multiplier} in [107] gives $2^am\in\mathcal D$
provided $H,G$ are nowhere zero,
\[
 4\max(\|H\|_\infty,\|G\|_\infty)+1<2^{a+1},\qquad
 \sum_{u\in U_a}w(G(u))=(m-1)\sum_{u\in U_a}w(H(u)).
\]
This result supplies the sign function explicitly. It takes the dyadic
sign functions $A,B$ associated with the signed binary digits of $1+4H$ and
$-1+4G$, respectively, and sets $f(n)=A(n)$ if $m\nmid n$ and
$f(n)=B(n/m)$ if $m\mid n$. The means of $A,B$ vanish by antisymmetry.
The valuation identity cancels the logarithmic layer statistic of
$mA+B-A(m\,\cdot)$, whose dyadic value function is the constant $m-2$.
The constant-value lemma in [107] then gives
$m\sum_{n\ge1}f(n)/n=0$.

\paragraph{Theorem.}
For each odd $m\ge3$ there is an explicit period $N=2^a m\in\mathcal D$
with $a\le2\lceil\log_2m\rceil+2$, and hence $N<16m^3$.
For $m=127$, the exponent $a=15$ suffices.

\paragraph{Proof.}
Put $v=\lfloor\log_2(m-1)\rfloor\ge1$, $J=v+2$, and $M=2^J$.
Multiplication by $m$ permutes $U_J$ in even cycles: the group is
a $2$-group and $1<m<M$. Negation pairs distinct cycles.
For $m\equiv1\pmod4$ this follows by reduction modulo $4$.
For $m\equiv3\pmod4$, even powers are $1\pmod4$, and for odd $b$
one has $v_2(m^b+1)=v_2(m+1)<J$, since $(m^b+1)/(m+1)$ is odd and
$m+1\le2^{v+1}$. Thus $-1$ never belongs to $\langle m\rangle$.

Select one cycle from each negative pair. Write a selected cycle
as $u_i=m^iu_0$, and group consecutive edges into pairs starting
at the even indices. Over all selected cycles there are $2^{v-1}$
such pairs. Give precisely $(m-1-2^v)/2$ pairs the index $t=v+1$,
and give the rest $t=v$. These are nonnegative integer counts.
The mean of $2^t$ over the pairs is exactly $m-1$.

On a pair $(u_{2j},u_{2j+1},u_{2j+2})$ set
$r_{2j}=1$, $r_{2j+1}=1+2^t$. Define $h(u_i)$ as the centered
representative of $u_i r_i\bmod M$, and extend by $h(-u)=-h(u)$.
Every value is odd, $|h|\le2^{v+1}-1$, and $h(u)\equiv u\pmod{2^v}$.
For either edge of the pair,
\[
 h(mu)-m h(u)\equiv mu(r_{i+1}-r_i)\pmod M.
\]
Its valuation is exactly $t$, including the edge closing the cycle,
because $t<J$. Negation preserves that valuation.
Set $H=2h$ and lift to $U_a$ for $a\ge J$ by reduction. This function is
$1$-Lipschitz in the stated congruence sense. For $j\le v$ this follows
from $h(u)\equiv u\pmod{2^v}$. For $j=v+1$, the factor $2$ suffices.
For $j\ge J$ the two values coincide. Each $w(H)$ equals $2$, and on
the two edges with index $t$ each $w(G)$ equals $2^{t+1}$. The pair
average proves the valuation identity, which is preserved under lifting.

Write $\ell=\lceil\log_2 m\rceil=v+1$. Then
\[
 \max(\|G\|_\infty,\|H\|_\infty)
 \le 2(m+1)(2^\ell-1)
 \le 2^{\ell+1}(2^\ell-1).
\]
Here $m+1\le2^\ell$ since $m\ge3$ is odd.
Consequently $a=2\ell+2$ satisfies the strict size inequality.
Since $2^\ell<2m$ the asserted period bound follows.

For $m=127$, put $M=256$. Start with
$H(u)=2\chi_4(u)$ on odd residues, where $\chi_4(1)=1$,
$\chi_4(3)=-1$ modulo $4$, and change just $H(127)$ to $126$ and
$H(129)$ to $-126$. The perturbations are multiples of $128$, so
$H$ remains $1$-Lipschitz. It remains antisymmetric and $w(H)=2$
everywhere. Since $127^2\equiv1\pmod{256}$, only the four inputs
$1,127,255,129$ have changed $G$-values, namely
$-128,-16000,128,16000$, respectively. Their weights are $128$ and
the other $124$ weights are $256$. Hence
$4\cdot128+124\cdot256=126\cdot128\cdot2$.
The maximum absolute value is $16000$ and $64001<2^{16}$,
so $127\cdot2^{15}=4\,161\,536$ is an admissible period.
The least exponent at $127$ remains undetermined here.

\paragraph{Optimality within the transfer criterion.}
For $m=1+3\cdot2^k$, $k\ge1$, any transfer certificate as above has
$\|H\|_\infty\ge2^{k-1}$. In fact $mu\equiv u\pmod{2^k}$.
If $a<k$, the two inputs coincide already in $U_a$. If $a\ge k$,
the $1$-Lipschitz condition and $2\|H\|_\infty<2^k$ force equal
integer values. In both cases $H(mu)=H(u)$ throughout, which would
give $w(G)=2^k w(H)$ and contradict the required multiplier
$m-1=3\cdot2^k$. At a point where $|H|$ is maximal,
$|G|\ge(m-1)\|H\|_\infty\ge(m-1)^2/6$. Thus
$a>2\log_2(m-1)-\log_2 3$. The leading constant $2$ in the universal
depth bound is sharp for this transfer method over odd multipliers.
The unrestricted least admissible period and its depth for prime
multipliers are separate questions.

\paragraph{Exact minimum within the criterion for Mersenne multipliers.}
For every $k\ge3$ put $q=2^k$ and $m=q-1$.
The least depth of a signed-digit transfer certificate of [107] is
\[
 a_{\rm cert}(2^k-1)=2k+1 .
\]
For the upper bound, on $U_{k+1}$ begin with $H=2\chi_4$, change
$H(m)$ to $q-2$ and $H(-m)$ to $-(q-2)$, and lift by reduction.
Both perturbations are multiples of $q$, which preserves the
$1$-Lipschitz condition. All values have weight $2$.
The four exceptional inputs $\{\pm1,\pm m\}$ give
$G(1)=-q$, $G(m)=-q(q-3)$ and their negatives. On the remaining
$q-4$ inputs, the weight of $G$ is $2q$. The weight mean is
$[4q+(q-4)2q]/q=2(q-2)$, as required. The maximum absolute value
$q(q-3)$ satisfies $4q(q-3)+1<4q^2=2^{2k+2}$, so $a=2k+1$ works.

For the converse, a certificate of smaller depth could be
lifted to depth $a=2k$, so assume this depth. The size bound gives
$\|G\|_\infty\le q^2/2-1$. The maximum principle
$\|G\|_\infty\ge(q-2)\|H\|_\infty$ yields
$\|H\|_\infty\le q/2+1$. The congruence $mu\equiv-u\pmod q$,
antisymmetry, and the Lipschitz property give $q\mid G(u)$
for every unit $u$. All values of $H$ have the same parity.
If that parity were odd, the required average of $w(G)$ would be
$q-2$, strictly smaller than its termwise lower bound $q$.
Thus $H$ is even-valued and $\|H\|_\infty\le q/2$.
It factors through $U_{k+1}$ since values above a fixed residue
modulo $2q$ are congruent modulo $2q$ and differ by at most $q$.

On $U_{k+1}$ one has $q\mid H(mu)+H(u)$ and
$|H(mu)+H(u)|\le q$. If $H(mu)=-H(u)$ then $w(G(u))=q w(H(u))$.
Otherwise $H(mu)=H(u)=\pm q/2$ and $w(G(u))=q$,
a loss of $q(q/2-1)$ from $q w(H(u))$. These exceptional inputs form
whole four-element orbits under multiplication by $m$ and negation:
$m^2\equiv1\pmod{2q}$ and $m\not\equiv\pm1\pmod{2q}$.
Write $E$ for their number. The budget equation now gives
\[
 0<E=\frac{2\sum_{u\in U_{k+1}}w(H(u))}{q(q/2-1)}
 \le\frac{2q}{q-2}<4 .
\]
This contradicts $4\mid E$. The minimum here concerns the transfer
criterion; an arbitrary vanishing periodic sign sequence need not
arise from that criterion.



\section{A density bound for the proposed minimal periods \(12p\)}
For a prime $p\ge5$, put $U=\mathbb F_p^\times$ and
$S=\langle2,3\rangle\le U$. We employ the period set $\mathcal D$ of the
preceding section.

\paragraph{Theorem.}
If $|S|$ is even then $12p$ is a divisibility-minimal member of
$\mathcal D$. The primes satisfying this hypothesis have natural density
$195/224$ among all primes. Consequently the set of primes satisfying
$12p\in\mathcal D$ has lower natural density at least $195/224$.

\paragraph{Proof.}
The cyclic group $S$ has a nontrivial quadratic character
$\eta:S\to\{-1,1\}$. Set $\alpha=\eta(2)$ and $\beta=\eta(3)$.
Since $2,3$ generate $S$, $(\alpha,\beta)\ne(1,1)$.
Choose one representative $c$ of each coset of $S$ in $U$ and define
$\lambda(ch)=\eta(h)$. This sign function satisfies exactly
\[
 \lambda(2y)=\alpha\lambda(y),\qquad
 \lambda(3y)=\beta\lambda(y),\qquad
 \sum_{y\in U}\lambda(y)=0.
\]
The last equality holds on each individual coset. In particular both
signs occur. A character of the entire group $U$ is unnecessary.

Apply the matching column of the three certificates for $12p$ in item
[093], Section~7 (the table entitled \emph{The three certificates for
$N=12p$}). Here is why substitution of this $\lambda$ is legitimate.
There the six functions on $U$, indexed by $d\mid12$, are
$\Psi_d(y)=\epsilon_d\lambda(y)(1-2{\boldsymbol1}_{P_d}(y))$,
where $P_d=\{2^i3^jy_0:(i,j)\in B_d\}$ and $\lambda(y_0)$ is the chosen
sign in the column. The constant-coefficient part of the first reduction
equation is zero by the displayed eigenvalue identities for translations.
The finite corrections satisfy the same Laurent polynomial identity
because translating each corrected term by $2^i3^j$ multiplies its
$\lambda$ factor by $\alpha^i\beta^j$. The total in layer $d$ is
$-2\epsilon_d\lambda(y_0)\sum_{(i,j)\in B_d}\alpha^i\beta^j$,
as the unrestricted $\lambda$-sum vanishes. Hence the remaining
mean and logarithmic equations coincide with the corresponding column.
All required distinctness checks for the points in a layer hold for
$p\ge7$ by the integer differences in that table. For $p=5$ one has
$(\alpha,\beta)=(-1,-1)$, for which the distinctness check also holds.
The reduction lemma of [093] produces a vanishing $12p$-periodic
sign function. By [093], its Theorems A and B and the exclusion of
periods congruent to $2$ modulo $4$, the only possible proper divisors
$4,12,4p$ do not belong to $\mathcal D$. This proves minimality.

The density of $|S|$ even is the published value $195/224$ in
F.~Pappalardi, \emph{Divisibility of reduction in groups of rational
numbers}, Introduction. The six reduced classes modulo $24$ covered by
[093] have total prime density $3/4$ and all lie in the present criterion,
because at least one of $2,3$ is a quadratic nonresidue. The additional
density is $27/224$. The criterion leaves the odd-order subgroup case
undecided; it yields no converse.



\section{Two Erd\H{o}s conjectures on abelian Cayley graphs}
For a simple graph $G$ of order $n$, write $b(G)$ for the minimum number of
edges whose deletion makes $G$ bipartite, and $g_G(k)$ for the minimum
number of induced edges among its $k$-vertex sets. Problems
\href{https://www.erdosproblems.com/23}{\#23} and
\href{https://www.erdosproblems.com/128}{\#128} ask, respectively, whether
$b(G)\le n^2/25$ and $g_G(\lfloor n/2\rfloor)\le n^2/50$ for every
triangle-free $G$. For a finite abelian group $A$ and
$S=-S\subseteq A\setminus\{0\}$, the graph $\operatorname{Cay}(A,S)$ joins
$x$ to $x+s$ for $s\in S$. It is $d$-regular with $d=|S|$, and it is
triangle-free exactly when $(S+S)\cap S=\varnothing$.

\paragraph{Theorem A (exponent five).}
Let $G=\operatorname{Cay}(\mathbb F_5^r,S)$ be triangle-free, $r\ge1$,
$n=5h=5^r$, and $m=nd/2$. Then $b(G)\le n^2/25$, with equality attainable
in every dimension. If $d>6n/25$, then
\[
 b(G)=m/5,\qquad
 g_G(k)=\frac d2\max\{0,k-2h,2k-n\}\quad (0\le k\le n).
\]
In particular, in this density range the optimal half spans exactly
$d(h-1)/4\le n^2/50$ edges.

\noindent\emph{Proof.}
For each pair $\{s,-s\}\subset S$, its edges split into $h$ disjoint
5-cycles, one on each coset of $\langle s\rangle$. These $d/2$ systems
partition $E(G)$. Every cut leaves at least one edge per cycle uncut,
giving $b(G)\ge m/5$. On a single 5-cycle the minimum numbers of induced
edges at cardinalities $j=0,\ldots,5$ are $0,0,0,1,3,5$. Each is at least
both $j-2$ and $2j-5$. Summing over a system gives the three lower bounds
$0,k-2h,2k-n$, and summing over directions gives the stated lower bound for
$g_G(k)$.

If $d>6n/25$, Versteegen's
\href{https://doi.org/10.1093/qmath/haae042}{Theorem 1.4 (2024)}
puts $S$ in the union of two cosets of a hyperplane. Since $S=-S$ and
$d>h$, its image in $\mathbb F_5$ is $\{a,-a\}$ with $a\ne0$.
Rescale the quotient so $a=1$. Its five fibers $V_i$ each have size $h$.
Each vertex has $d/2$ neighbors in each adjacent fiber and no others.
The five consecutive fiber pairs carry $m/5$ edges each. Deleting one
such pair leaves a graph mapping to a path, hence $b(G)=m/5$. We have
$d\le2h$, so $m/5\le n^2/25$, with equality for the full preimage of
$\{1,-1\}$. For $k\le2h$, choose vertices in $V_0\cup V_2$. For
$2h\le k\le3h$, take $V_0\cup V_2$ and $k-2h$ vertices of $V_4$. Each
added vertex contributes $d/2$ edges. For $k\ge3h$, take the complement
of an independent set of size $n-k$ in $V_0\cup V_2$. Its induced edge
count is $m-(n-k)d$. These choices meet all lower bounds.

For $d\le6n/25$, replace each vertex by an independent set of $t$ twins
and each edge by a complete bipartite graph. The resulting triangle-free
graph $G[t]$ has $N=tn$ vertices and edge density
$e(G[t])/\binom{N}{2}=(d/n)N/(N-1)$. As $6/25<0.2486$, for all
sufficiently large $t$ this density is at most $0.2486$.
\href{https://arxiv.org/abs/2103.14179}{Balogh, Clemen and Lidick\'y,
Theorem 2(c) (2021)} give $b(G[t])\le N^2/25$. We have
$b(G[t])=t^2b(G)$: lifting a cut gives one inequality, and selecting
one uniformly random representative from each fiber in an optimal
blow-up cut gives the other by expectation. Dividing by $t^2$ yields
the claim.\hfill$\square$

\paragraph{Theorem B (exponent dividing four).}
Let $G=\operatorname{Cay}(A,S)$ be triangle-free, where $A$ is a nontrivial
finite abelian group of exponent dividing four. Write $n=|A|$ and let
$q$ be the least eigenvalue of the signless Laplacian $Q=dI+A_G$.
There is a balanced optimal cut, and
\[
 b(G)=nq/4,\qquad g_G(n/2)=nq/8.
\]
Consequently $b(G)\le15n^2/376<n^2/25$ and
$g_G(n/2)\le15n^2/752<n^2/50$.

\noindent\emph{Proof.}
The edgeless case is immediate. The characters diagonalize $Q$ with
real eigenvalues $d+\sum_{s\in S}\chi(s)$. A minimizing character is
nontrivial: the trivial character has eigenvalue $2d$, above the mean
eigenvalue $d$. For a real minimizing character take $f=\chi$. For a
character of order four take $f=\sqrt2\operatorname{Re}(e^{\pi i/4}\chi)$.
Its four possible phases give only the values $1,-1$. In either case
$Qf=qf$ and $\sum_x f(x)=0$. For any sign vector $u$, the number of uncut
edges is $u^{\mathsf T}Qu/4\ge nq/4$ because
$u^{\mathsf T}Qu=\sum_{xy\in E}(u_x+u_y)^2$. The cut given by $f$
attains equality. The two sides of every balanced cut have equal
induced edge counts by regularity, proving the formula for $g_G(n/2)$.
The published bound $q\le15n/94$ follows from
\href{https://doi.org/10.1137/22M150767X}{Balogh, Clemen, Lidick\'y,
Norin and Volec (2023)}.\hfill$\square$

\paragraph{Corollary (normal Cayley graphs).}
Theorem B holds for any finite nontrivial group $K$ of exponent dividing
four if $S=S^{-1}\subseteq K\setminus\{1\}$ is invariant under
conjugation and its Cayley graph is triangle-free.

\noindent\emph{Proof.}
Assume $d>0$ and use right translations $R_xF(g)=F(gx)$. The adjacency
operator $\sum_{s\in S}R_s$ comes from a central element of the group
algebra. By Schur's lemma it acts as a scalar on each irreducible
representation. Pick one on which $Q$ acts by $q$. The group $K$ is a finite
$2$-group by Cauchy's theorem, hence supersolvable. Each irreducible
complex representation of a finite supersolvable group is induced
from a linear character of a subgroup
(\href{https://www.uni-math.gwdg.de/tammo/rep.pdf}{T. tom Dieck,
\emph{Representation Theory}, Theorem 4.3.1}). Realize the chosen
representation as the right-translation space of functions
$F(hg)=\psi(h)F(g)$ for $h\in H$, with $\psi$ a linear character of
$H\le K$. Specify $F(t)=1$ on representatives $t$ of $H\backslash K$.
All values of $F$ are fourth roots of unity since $h^4=1$. The
real vector $f=\sqrt2\operatorname{Re}(e^{\pi i/4}F)$ has sign
entries and $Qf=qf$. Since $q\le d<2d$ and $Q\mathbf 1=2d\mathbf 1$,
symmetry makes $f$ orthogonal to $\mathbf 1$. The cut and half identities
now follow as in Theorem B.\hfill$\square$

The October 2026 preprint of
\href{https://arxiv.org/abs/2610.14503}{Cai, Chen and Zhang}
gives $q\le7n/50$ for every triangle-free graph. With that input the
two bounds in Theorem B become $7n^2/200$ and $7n^2/400$. In particular,
its balanced cut continues to certify both Erd\H{o}s inequalities after
arbitrary edge deletions and at most $\lfloor n^2/200\rfloor$ edge
additions: the total within the two sides is then at most $n^2/25$, and
one side contains at most half that total. Using only the 2023 bound
gives the same conclusion with $\lfloor n^2/9400\rfloor$ additions.

\clearpage

\section{Pinned distances on a conic and in convex octagons}
For a finite set $X$ in the plane, the pinned distances at $P\in X$ are the distinct numbers $|P-Q|$ for $Q\in X\setminus\{P\}$. Erd\H{o}s problem \href{https://www.erdosproblems.com/982}{982} asks for at least $\lfloor n/2\rfloor$ pinned distances at some vertex of every convex $n$-gon.
\paragraph{Theorem.}
This bound holds for every strictly convex polygon whose vertices lie on one nonsingular real conic. If $n\ge5$, one can choose the vertex so that each positive distance occurs at most twice there.
\paragraph{Proof.}
For an ellipse use coordinates $(au,\pm b\sqrt{1-u^2})$, $a\ge b>0$. Choose a point with the largest $|u|$ in $X$. By reflections write it $P=(ac,bs)$ with $c,s\ge0$, $s=\sqrt{1-c^2}$. All points lie in the strip $-c\le u\le c$. Along either circular semicircle inside this strip, the squared chord distance from $(c,s)$ is a strictly decreasing function of $u$ when $c>0$. Indeed its derivative on the interior is $-2c\pm 2su/\sqrt{1-u^2}<0$ there, and for $c=1$ the derivative is $-2$. Elliptic squared distance adds the decreasing term $(a^2-b^2)(u-c)^2$ to $b^2$ times the circular squared distance. A circle centered at $P$ therefore encounters at most one sampled point on each semicircle. The case $c=0$ has at most two points.
  
For a parabola or a single branch of a hyperbola, use orthogonal coordinates in which the curve is $y=h(|x|)$ with $h$ increasing. Choose the point $P=(t_0,h(t_0))$ of least $|x|$, reflecting the horizontal axis coordinate if necessary. For $t\ge t_0$, the squared distance to $(\pm t,h(t))$ is $(t\mp t_0)^2+(h(t)-h(t_0))^2$, strictly increasing on each side.
A strictly convex set with at least five vertices cannot use both branches of a hyperbola. For this incidence claim only, an affine transformation sends the hyperbola to $P(t)=(t,1/t)$. If $0<u<v<w$ and $z<0$, then $P(v)$ is strictly inside the triangle with vertices $P(u),P(w),P(z)$. This follows by testing the three oriented edges with
\[
 \det(P(t)-P(s),P(r)-P(s))
 =\frac{(t-s)(r-s)(r-t)}{str}.
\]
The case of three negative parameters follows by central reflection. Finally, for $n\le4$ the bound is immediate: four distinct planar points cannot all have the same pairwise distance.\hfill$\square$
\paragraph{The general octagon.}
Every convex octagon has a vertex with at least four pinned distances, whether or not it lies on a conic. By the Erd\H{o}s--Fishburn and Fishburn classification, as stated in \href{https://arxiv.org/abs/0906.0199}{Nozaki and Shinohara, Proposition 4.18}, all eight-point planar sets with at most three distances at each vertex are similar to one configuration. To see its failure of convex position explicitly, put $t=1+\sqrt3$ and take the four points $(\pm1,\pm1)$ and the four axis points $(\pm t,0),(0,\pm t)$. At each inner point the positive distances are $2,2\sqrt2,\sqrt2 t$, and at each outer point they are $2,\sqrt2 t,2t$. This realizes the classified configuration. The four inner points lie strictly inside the outer diamond, since $2<t$. The regular octagon attains four at each vertex, so the octagon minimum is exactly four. This is a consequence of the existing classification. The general convex-polygon question remains open.

\clearpage
\appendix
\section{Exact residue verification for the consecutive-sum constructions}
The following Python program checks interval intersections and modular sums with exact rational arithmetic. It verifies all 35 nonempty cases for the stronger table. The bounded-deletion argument in Section 2 turns the interior checks into the all-$N$ statement.
{\footnotesize\begin{verbatim}
"""Verify both lower bounds by rational interval arithmetic and periodic sums."""
from fractions import Fraction as F
from math import lcm


def mask(q, residues=()):
    return q, frozenset(residues)


def complement(pattern):
    q, residues = pattern
    return mask(q, set(range(q)) - residues)


def window_sums(pattern, length, target_modulus=None):
    q, residues = pattern
    m = target_modulus or length * q
    period = lcm(q, m)
    count = period // q * len(residues)
    terms = [n for n in range(period * (length + 1)) if n % q in residues]
    return mask(m, {sum(terms[i:i + length]) % m for i in range(count)})


def excluded_sum_union(first, k, second, j):
    m = lcm(k * first[0], j * second[0])
    forbidden = window_sums(first, k, m)[1] | window_sums(second, j, m)[1]
    return complement(mask(m, forbidden))


base = [
    (F(1, 5), mask(1)),
    (F(2, 9), mask(16, [0, 1, 2, 5, 6, 8, 10, 11, 14, 15])),
    (F(1, 4), mask(1, [0])),
    (F(8, 27), mask(1)),
    (F(1, 3), mask(4, [1, 2, 3])),
    (F(3, 8), mask(3, [1, 2])),
    (F(4, 9), mask(32, [0, 2, 6, 8, 9, 12, 13, 16, 19, 20, 23, 24, 26, 30])),
    (F(1, 2), mask(2, [0])),
    (F(16, 27), mask(1, [0])),
    (F(2, 3), mask(8, [1, 2, 4, 6, 7])),
    (F(3, 4), mask(3, [1, 2])),
    (F(8, 9), complement(mask(64, [2, 8, 14, 17, 21, 25, 29, 35, 39, 43, 47, 50, 56, 62]))),
    (F(1), mask(4, [0, 1, 3])),
]


def modify(rows, interval, scale, pattern):
    low, high = interval
    rows.append((scale * low, scale * high, pattern))


def verify(patches, expected):
    endpoints = {F(0), *(e for e, _ in base)}
    for low, high, _ in patches:
        endpoints.update([low, high])
    for i, (low, high, _) in enumerate(patches):
        assert all(max(low, a) >= min(high, b) for a, b, _ in patches[:i])

    def lookup(x):
        for low, high, pattern in patches:
            if low < x < high:
                return pattern
        return next((p for end, p in base if x < end), mask(1))

    density = F(0)
    cuts = sorted(endpoints)
    for low, high in zip(cuts, cuts[1:]):
        q, residues = lookup((low + high) / 2)
        density += (high - low) * len(residues) / q
    assert density == expected
    # Every selected integer exceeds N/6, so higher block lengths cannot hit.
    assert all(not lookup((a + b) / 2)[1] for a, b in
               zip(cuts, cuts[1:]) if a < F(1, 6))
    checks = 0
    for k in range(2, 6):
        grid = sorted(endpoints | {e / k for e in endpoints})
        for low, high in zip(grid, grid[1:]):
            x = (low + high) / 2
            source = lookup(x)
            modulus, target = lookup(k * x)
            if source[1] and target:
                residues = window_sums(source, k, modulus)[1]
                assert residues.isdisjoint(target), (k, low, high)
                checks += 1
    print(f"{density}: {checks} nonempty source/target cases verified")


patches = []
I = F(16, 81), F(1, 5)
X = mask(16, [1, 4, 6, 9, 14])
U = mask(32, [0, 3, 6, 9, 11, 16, 17, 22, 24, 28])
for scale, pattern in [(1, X), (2, U), (4, complement(window_sums(U, 2)))]:
    modify(patches, I, scale, pattern)
verify(patches, F(703, 1296))

Y = mask(64, [29, 33, 37, 41])
J = F(3, 16), F(16, 81)
H = F(8, 45), F(3, 16)
K = F(1, 6), F(128, 729)
A, B = mask(4, [2]), mask(12, [2, 10])
Cp, Up, W = mask(192, [88, 152]), mask(8, [3, 5]), mask(8, [0])
for interval, source in [(J, A), (H, B)]:
    modify(patches, interval, 1, source)
    modify(patches, interval, F(3, 2), Y)
    modify(patches, interval, 3, excluded_sum_union(source, 3, Y, 2))
modify(patches, J, 2, Up)
modify(patches, J, 4, complement(window_sums(Up, 2)))
modify(patches, K, 1, Cp)
modify(patches, K, F(3, 2), W)
modify(patches, K, 3, complement(window_sums(W, 2)))
verify(patches, F(3040369, 5598720))
\end{verbatim}}

\end{document}
